
Deep learning models can achieve high overall accuracy while exhibiting poor performance on specific subgroups. This failure mode arises when training data contains spurious correlations---statistical associations between input features and labels that do not hold across all subpopulations. For example, in the Waterbirds benchmark, 95\% of waterbird images appear on water backgrounds, leading models to associate ''water background'' with ''waterbird'' rather than learning features of the birds themselves. Models exploiting this shortcut achieve high aggregate accuracy but fail on the minority of waterbirds appearing on land backgrounds.

Group Distributionally Robust Optimization (GroupDRO) addresses this problem by minimizing worst-group loss during training, but requires group annotations identifying which samples belong to minority versus majority groups. Without such annotations, practitioners lack training-time signals indicating whether their models are learning spurious correlations. The failure mode remains invisible until deployment on data where the spurious correlation does not hold.

This work investigates whether per-sample loss trajectories during standard training encode information distinguishing minority from majority samples. The underlying hypothesis is that minority samples---where spurious features conflict with true labels---experience different optimization dynamics than majority samples. If loss trajectory features can predict minority group membership, they could serve as a diagnostic tool for spurious correlation vulnerability.

The experiments reported here address three questions:

\begin{enumerate}
\item \textbf{Existence:} Do per-sample loss trajectory features predict minority group membership with discriminative power above chance?
\item \textbf{Feature analysis:} Which trajectory features are most informative?
\item \textbf{Spurious-specificity:} Is the trajectory signal specific to spurious correlation, or does it reflect generic sample difficulty?
\end{enumerate}

The main findings are as follows. First, trajectory features extracted from epochs 1-5 predict minority membership with AUROC = 0.9452 on Waterbirds, exceeding the pre-specified 0.75 threshold. Second, initial loss ($L_{1})$ alone achieves AUROC = 0.9473, suggesting the signal is present from the first training epoch. Third, the signal attenuates substantially under GroupDRO ($\Delta AUROC$ = 0.29) but minimally under variance-matched random reweighting ($\Delta AUROC$ = 0.01), providing evidence for spurious-specificity. Fourth, a secondary hypothesis about delayed curvature stabilization in minority samples was not supported by the data.
